\honest{Faster Than Science}{Eliezer Yudkowsky}{2008}

I sometimes say, I begin, that science piles up so much evidence that even scientists cannot ignore it. Max Planck was gloomier: new truths win because their opponents die. I am ``much tickled'' by this, since it means that science rests ``on the good taste of grad students.'' I quote Planck without a source, as if it were a finding. Studies that tested his remark found mixed results.

Physicists, I say, are gradually coming to accept the many-worlds interpretation of quantum mechanics, which ``suggests'' that some of them need a large enough pack before they will convert. I give no survey and no numbers. Seventeen years later, a survey of more than 1,100 physicists found 15\% favoring many-worlds, well behind the older Copenhagen interpretation.

Scientists are moved by prejudice and herd behavior, I say, and each unjustified shift must be undone with extra evidence, so academic science is ``a highly inefficient processor of evidence.'' Inefficient compared with nothing that exists: my standard is an ideal reasoner. The collapse of the wavefunction, I add, has no experimental support, and one argument against it is that it violates relativity. A version of collapse consistent with relativity had been published two years earlier.

Then an accusation. In some fields debates drag on for decades because ``one side keeps throwing up more and more ridiculous objections'' from ``an entrenched position of academic power.'' I am thinking of evolutionary psychology, I say. I name no one, no paper and no objection.

Can individuals reach the truth faster than science? My answer is about where hypotheses come from. Science tests ideas but does not say where they come from. That is true, and philosophers have discussed it since the 1930s.

My one real idea follows. When the possible answers are many, finding the right one takes more evidence than confirming it. Among four billion equations, getting one of them up to 10\% takes far more evidence than taking it from 10\% to 90\%. This is correct, and Charles Peirce made the same point in 1903. It also assumes that the scientist starts with every equation equally likely. No scientist does. They start from what their training makes plausible, and the training comes from the slow public process I say they outrun.

Suppose, I say, that a robot-controlled Ouija board kept producing hypotheses that passed their tests. Science would not blink, but ``The pure ideal essence of Bayes would burst into flames and die.'' A Bayesian would simply conclude that the Ouija board was reliable. ``Compared to Science,'' I add, ``Bayes is falsified by more of the possible outcomes.'' Bayes's rule is a theorem, and no outcome falsifies it.

At the frontier, I say, scientists see things that are not yet confirmed. A scientist who holds the right hypothesis at 10\% ``must, in fact, have done'' most of the work, or he would never have thought of it, though he ``may not know'' that he has. So in the interval before the experiment, he ``rationally knew'' something science did not. At 10\%, he expects to be wrong nine times in ten. Many scientists reason the same way and guess wrong, and I count only the one who guessed right. His guess becomes knowledge when he ``performs the experiment, publishes the result.''

I end by limiting all this to non-routine science. ``It is much easier to train people to test ideas, than to have good ideas to test.'' That is true.
